% Selected-scale charts and the public gallery handoff.
\section{From funded chart receipts to a trace gallery}
\label{sec:ex-lattice-firesky-gallery}

This section records the actual selected-run route from one-scale chart data
to the scale-dependent charts used by the trace calibration. The analytic
chart receipts are provided by the funded construction and are identified
below. We prove the subsequent scale selection and finite-dimensional limit
steps explicitly. The correction and Ward-flux calculations are in
Section~\ref{sec:ex-lattice-trace-producer}.

\subsection{The funded chart family and its one-step output}

Fix a grade $d$, a binding, a gauge-invariant chart budget, and a scale
parameter $z\geq0$. The funded chart receipt consists of three panels
$P(t,i)$, $i=0,1,2$, for $0\leq t\leq1$, with arc coordinates $b_i+t$.
Their currency and angular parameters are constant in $t$, and the three
base arcs obey the finite separation inequalities for the probe count at
grade $d$. At every sufficiently large level $k$ and every $t\in[0,1]$,
the receipt supplies a finite-volume chart witness at that level and
offset. The other fields of this witness are the chart, invoice,
legacy-cost, actuation, and stamp data used by the quantitative chart
estimate. The single offset changes only the three arc coordinates.

The funded run constructs these panels from its paid multi-scale request.
The request gives charts with a common added offset and fixed currency and
angle. Forward chart estimates supply legacy costs at all large levels;
the staging estimate at the common angle supplies simultaneous actuation
and stamp receipts. Clipping the offset to $[0,1]$ preserves the arc lower
bound, chart factor, currency, and angle. The resulting finite-level data
are the hypotheses of the one-step chart theorem.

That theorem also requires a simple gauge group, the funded chart-control
estimate, a bound on the generator through grade $d$, an index $l$ with
$l\leq\operatorname{trays}(d)$, a positive sequence tending to zero, and
the selected extraction order, smooth table, point limits, and reference
hierarchy used by the station construction.
It returns a strictly increasing level curve and, at a selected point,
data $(k,i,b,c,R)$ such that the requested scale equals the exact clock
gap at level $i$ and base arc $b$. Its chart has $5<b$ and $b$ no larger
than the permitted angular bound, covers all reference grades through
$d$, and supplies the finite-volume charts and invoices whenever $i<N$.
It also supplies a pairwise comparison receipt at cost
$\operatorname{costHint}(l)$ and a station whose lift has a finite norm
bound, preserves the unit, real structure, and signed action, and obeys
the short-test estimate for arity at most $l$, tests in the stated window,
and error $\operatorname{costHint}(l)$. The proof of this quantitative
upgrade is the separate three-panel chart result; this section uses that
exact output and proves the scalar gap interpolation and recursion. In
the gallery construction one takes $l=d$, so the chart cutoff, arity cap,
and finite grade cutoff agree. The clock gives, for each fixed base arc
$b\geq0$, a positive gap
\begin{equation}\label{eq:ex-lattice-firesky-gap}
 g_b(k)=a_k m_k^b,
\end{equation}
where $a_k>0$, $m_k\geq1$, $g_b(k)\to0$, and eventually
$g_b(k)\leq g_{b+1}(k+1)$. These scale-gap facts are the input to the
selection argument.

\begin{lemma}[Clock-gap interpolation]
\label{lem:ex-lattice-firesky-gap-interpolation}
Fix $b\geq0$. Suppose $g_b(k)>0$, $g_b(k)\to0$, and
$g_b(k)\leq g_{b+1}(k+1)$ for all $k\geq k_0$. For every $J\geq k_0$
there is $r>0$ such that for each $0<s\leq r$ there are $k>J$ and
$t\in[0,1]$ with $s=g_{b+t}(k)$.
\end{lemma}
\begin{proof}
Take $r=g_b(J)>0$. Given $0<s\leq r$, choose the least $k>J$ with
$g_b(k)\leq s$; it exists by convergence to zero. Minimality gives
$s\leq g_b(k-1)$, and the one-step comparison at $k-1$ gives
\[
 g_b(k)\leq s\leq g_b(k-1)\leq g_{b+1}(k).
\]
The function $t\mapsto g_{b+t}(k)=a_km_k^{b+t}$ is continuous on
$[0,1]$. The intermediate value theorem supplies the desired $t$.
\end{proof}

\subsection{The fields of a chart witness}
\label{subsec:ex-lattice-firesky-witness}

The producers of the pairwise receipt use several fields of the
finite-volume chart witness. We list them here, with the two scale
definitions they require, and prove the one property that follows from
definitions alone. The witness itself is part of the funded chart receipt
above, which remains an input of this section. In this subsection, $k$
denotes the level of the chart and $\mu_k$ the monitor $m_k$ of
\eqref{eq:ex-lattice-firesky-gap}. $J$ is the stage clock and $L_j$ are
the block lengths of Section~\ref{sec:common-scale-selection}, and
$c_k\ge1$ is the cover size of the clock at level $k$.

\paragraph{Two levels attached to a chart.}
Let the chart have base arc $b$, and put $m=\mu_k^b$. Its \emph{middle
level} and \emph{root depth} are
\[
 \ell=\max\Bigl\{j\le k:\frac{a_j}{a_k}\ge m\mu_k^{2}\Bigr\},\qquad
 w=\max\Bigl\{j\le\ell:\frac{a_j}{a_\ell}\ge c_k^{10}\Bigr\},
\]
with the convention $\max\varnothing=0$. Thus $w\le\ell\le k$. The middle
cell of the chart is a cell of spacing $a_\ell$. The root tests and the
chamber condition of the chart are taken on the centered tiling at depth
$w$, whose cells are indexed by the sites of the level-$w$ torus.

\paragraph{The witness fields.}
A chart witness at level $k$ with base arc $b$ supplies the following
fields. Here $M$ is the marker allowance of the chart, $N>k$ and $p$ are
arbitrary.
\begin{enumerate}
\item[(F1)] \emph{Allowance.} Four conditions hold:
 \begin{itemize}
 \item dilation: $m\mu_k^2\le a_\ell/a_k$, so the set defining $\ell$
  is nonempty;
 \item depth: $c_k^{10}\le a_w/a_\ell$;
 \item wheel: the smaller of the clock factors of the stages $J(k)$ and
  $J(k)-1$ is at least $2$;
 \item bracket: $J(k)-1\le J(j)$ for every $w\le j\le k$.
 \end{itemize}
\item[(F2)] \emph{Plateau.} $\widehat\beta_k/2\le\widehat\beta_j\le\widehat\beta_k$
 for every $w\le j\le k$.
\item[(F3)] \emph{Lawful history.} For each $N>k$ and $p$ there are
 representations $P_j$ on the level-$j$ torus, $w\le j\le k$, along the
 course of the chart. Each $P_j$ satisfies the suppositions of
 Proposition~\ref{prop:rank-root-output-ticket}(2) at level $j$, with the
 clip and row bound of the stage $\pi_{J(j)}$. Its budget radius is
 $\lceil(\log\widehat\beta_j)^2\rceil^2$, and its threshold is the clock's
 hatted coupling at $j$ times the corresponding healthy allowance.
\item[(F4)] \emph{Ticket.} $T_M(\pi_{J(w)})\le\log\mu_k/(16M)$.
\item[(F5)] \emph{Readings and stamps.} For every grade $d'\le d$ there is
 a linear reading $T_{d'}$ of the grade-$d'$ label space with three
 properties:
 \begin{itemize}
 \item it fixes the reference unit;
 \item it commutes with the real structure;
 \item it satisfies the half-accuracy bound
  $\|a(v)-T_{d'}v\|\le\frac12\|a(v)\|$ in coefficient size, where $a$
  is the calibrated entry map at level $k$ for $(N,p)$.
 \end{itemize}
 It also carries a stamp duty with the same reading.
\end{enumerate}
By (F1), every level of $[w,k]$ lies in the stage $J(k)$ or $J(k)-1$. Hence
$J(w)\ge J(k)-1$, and $w\to\infty$ as $k\to\infty$ because $J$ diverges.

\paragraph{Corrected labels.}
Write $\sigma=\sigma_{k,N,p}=(\beta_{N,p,k}^{1/2}\,\theta_k)^{-1}>0$ for
the shrink factor, where $\theta_k$ is the clock threshold. Write $W_\rho$
for the linear map multiplying the grade-$j$ component of a label by
$\rho^j$, and $S_\sigma$ for the frame scaling by $\sigma$. The
\emph{corrected label} of a basis vector $e$ is
\[
 X_{N,p}e=T_{d'}^{-1}\bigl(S_{\sigma}^{-1}W_me\bigr).
\]
These are the linear maps $X^{(i)}_{N,p}$ of input (G2) in
Subsection~\ref{subsec:trace-rounded-comparison}, for the chart selected at index
$i$.

\begin{lemma}[Corrected labels solve the reading equation]
\label{lem:ex-lattice-firesky-solves}
The reading $T_{d'}$ is bijective, and $X_{N,p}e$ is the unique label $v$
with
\begin{equation}\label{eq:ex-lattice-firesky-solves}
 S_\sigma\bigl(T_{d'}v\bigr)=W_me .
\end{equation}
\end{lemma}
\begin{proof}
The calibrated entry map satisfies the reference inverse estimate
\eqref{eq:rank-reference-inverse} of
Section~\ref{sec:actual-rank-purchases}, $|v|_1\le D_k|a(v)|_1$. With the
half-accuracy bound of (F5), the argument given there shows that $T_{d'}$
is injective, hence bijective. Next, $S_\sigma$ is invertible. So the
equation has the unique solution $X_{N,p}e$; this is
\eqref{eq:rank-common-lifts} with $\rho=m$.
\end{proof}

\paragraph{Status.} Fields (F1)--(F5) are fields of the funded witness.
This section records them but does not construct them. Their source
construction is the cost and staging construction of the funded run,
which this section already takes as input. With these fields, three
of the inputs listed elsewhere are now displayed rather than missing:
\begin{itemize}
\item the allowance and the root-depth identification used in
 Subsection~\ref{subsec:trace-rounded-receipt};
\item the bracket property used in
 Subsection~\ref{subsec:rank-root-output-ticket};
\item the ticket at the output level.
\end{itemize}
The plateau (F2) and the lawful history (F3) are the properties that the
propagation to the output torus requires at every level of $[w,k]$.

\begin{lemma}[Chamber interiority at the root depth]
\label{lem:ex-lattice-firesky-chamber}
Consider a chart at level $k$ with base arc $b$, middle level $\ell$ and
root depth $w$. Let $x$ be fit for the chart: the grid configuration
$z=a_k\lceil mx\rceil$ is coordinatewise separated and
$|z_{c,\alpha}|\le a_\ell/8$ for $c=0,1$ and $\alpha=1,\dots,4$. Fix a grade
$d$, and let $Z_d$ bound the coordinates of the finite stencil $S_d$ of
the grade-$d$ insertions. Then, for every $N$ with
\begin{equation}\label{eq:ex-lattice-firesky-chamber-mesh}
 a_N(Z_d+2)<\tfrac38\,a_w ,
\end{equation}
every period $p$ and every grade-$d$ label, the insertion at the site
$z_c$ ($c=0,1$) has every link of its stencil, including both endpoints,
strictly inside the home chamber of the centered tiling at depth $w$. It
is therefore in the chamber-local domain of that tiling. Condition
\eqref{eq:ex-lattice-firesky-chamber-mesh} holds for all large $N$, with
a threshold depending only on $w$ and $d$. In particular (G4) of
Subsection~\ref{subsec:trace-rounded-comparison} holds with $w_i=w$ for
the chart selected at index $i$.
\end{lemma}
\begin{proof}
Measure positions in mesh steps of the level-$N$ lattice. The depth-$w$
cell has an odd number $q=a_w/a_N$ of steps along each axis, as in the
interiority argument of Section~\ref{ex-centered-construction}. Its home
cell consists of the offsets $u$ with $|u_\alpha|\le(q-1)/2$ from the
centre. Since $a_k/a_N$ is an integer, the site $z_c$ has integer
offset $u=z_c/a_N$. By fitness, and because $w\le\ell$ implies
$a_\ell\le a_w$,
\[
 |u_\alpha|\le\frac{a_\ell}{8a_N}\le\frac{q}{8}.
\]
A link of the stencil at $z_c$ has endpoints $u+s$ and $u+s+e$, with
$s\in S_d$ and $e$ a unit vector. Each coordinate of an endpoint is at
most $q/8+Z_d+1$ in absolute value. By
\eqref{eq:ex-lattice-firesky-chamber-mesh},
$Z_d+2<3q/8$. Hence
\[
 \frac q8+Z_d+1<\frac q2-1<\frac{q-1}2 ,
\]
so both endpoints lie strictly inside the home cell. The insertion is a
bounded, measurable, gauge-invariant function of these links. It is
therefore chamber-admissible, exactly as in the argument cited above.
Finally, $a_N\to0$ while $a_w>0$ and $Z_d$ are fixed, so
\eqref{eq:ex-lattice-firesky-chamber-mesh} holds for all large $N$. The
threshold does not depend on $p$, on the label, or on the fit
configuration.
\end{proof}

\subsection{Choosing scales and chart levels}

\begin{lemma}[Strict scale selection]
\label{lem:ex-lattice-firesky-selection}
Assume the one-step chart producer has this output at every grade $d$: for
every previous index $N$ and every $\theta>0$, it supplies an index $n>N$,
a level $i\geq d$, and a grade-$d$ chart with exact radius $s_n<\theta$
and the finite-level receipts above. Let $s_n>0$ and $s_n\to0$. Then
there are strictly increasing indices $n_d$, levels $i_d\geq d$, and
admissible charts at grade $d$ and level $i_d$, such that
\[
 s_{n_0}<1,\qquad s_{n_{d+1}}<s_{n_d},\qquad i_d\geq d.
\]
In particular $i_d\to\infty$ and the selected chart radius is exactly
$s_{n_d}$.
\end{lemma}
\begin{proof}
For $d=0$, apply the one-step output with previous index $0$ and threshold
$1$. For $d>0$, apply it with previous index $n_{d-1}$ and threshold
$s_{n_{d-1}}$. It gives $n_0>0$, $s_{n_0}<1$, and thereafter
$n_d>n_{d-1}$, $s_{n_d}<s_{n_{d-1}}$, with a grade-$d$ chart at level
$i_d\geq d$ and exact radius $s_{n_d}$. Induction proves the claims;
$i_d\geq d$ implies $i_d\to\infty$.
\end{proof}

The one-step chart theorem obtains this output by combining the clock-gap
interpolation with the funded chart family. Its quantitative comparison
and station estimates are the three-panel weak-comparison result in
Section~\ref{sec:three-panel-weak-comparison}; we use its conclusion here
without duplicating those estimates.

Here is the finite-level extraction used to pass from the charts to the
limiting fields. Let $F$ be an extraction order, let $\mathcal U$ be an
ultrafilter on period-phase pairs finer than the fused filter of $F$, and
fix a cutoff $D$. Suppose $X_{N,p}$ is a tuple of linear maps, one for
each grade $0,\ldots,D$, from the corresponding reference field into the
common finite-period coefficient space. Assume $A\geq0$ and, on an
$F$-late set,
\[
 \|X_{N,p}\|\leq A.
 \tag{1}\label{eq:ex-lattice-firesky-lift-bound}
\]
The finite product of these map spaces is finite-dimensional. The closed
ball of radius $A$ is compact. Compactness gives a proper subfilter limit
of the bounded tuple of maps. This subfilter is finer than $\mathcal U$;
ultrafilter maximality identifies it with $\mathcal U$. Thus the tuple
converges along $\mathcal U$ to $Y=(Y_0,\ldots,Y_D)$. This is the only
compactness step: it is finite-dimensional and is performed at each fixed
cutoff.

If the finite-level maps preserve the real involution and signed actions on
the late set, continuity of those actions shows that each $Y_d$ preserves
them. If they send the reference unit to the scalar unit eventually, the
same limit identity gives $Y_d(1)=1$. Compose with the quotient projection
to define the table map
\[
 U_d=\mathsf P_dY_d,\qquad 0\leq d\leq D.
\]
This order matters: the finite maps converge in the coefficient space, and
the table-valued maps are obtained by applying the continuous projection
after taking the limit.

For a fixed tuple of $n$ labels of grades at most $D$, a fixed smooth test
function supported in the allowed window, and $n$ no greater than the
available arity cap, the finite-level receipt bounds the paid test error
by $\epsilon$ times the label charge and the test supremum. To pass this
estimate to the quotient table, use the ordered-moment limit rather than
continuity of an integral whose finite-volume measure varies. For each
fixed tuple of coefficient vectors, the ordered read receipt gives
convergence of the corresponding finite-volume moment to the table moment.
The moment is multilinear in its $n$ label vectors, and each label space is
finite-dimensional. Expanding each moving vector in a fixed basis expresses
the moment as a finite sum of products of convergent coordinates and
convergent fixed-basis moments. Hence the moving-label moment converges to
the table moment. This is the finite-basis multilinear limit argument; it
uses the common ordered extraction and the fixed compactly supported test.
The inequality passes to the limit
because a closed interval is closed under limits. More explicitly, for
each $\eta>0$ the late-set estimate is at most the claimed bound plus
$\eta$; taking the limit and then letting $\eta\downarrow0$ removes the
slack. This proves the cutoff-level short-test estimate with no exchange
of an infinite sum and a limit. It is the same moment-passage mechanism as
the calibrated-smearing lemma: fixed-vector convergence is supplied by the
ordered table, and finite dimensionality handles the moving vectors.

Apply this extraction separately at cutoff $D=d$ to the chart chosen at
grade $d$ and level $i_d$. The chart receipt supplies a bound on its finite lifts,
exact reality and signed covariance, the exact unit, and a short-test error
for arity $n\leq d$ with coefficient $1/(d+1)$. Projecting the
limits gives panels $U_d$. For every fixed finite label tuple, its grades
are at most $d$ eventually along the sequence of cutoffs; the support of a
fixed compact test fits the expanding chart window eventually; and
$1/(d+1)\to0$. Write
$\mathcal T_{r,U}(v;t)$ for the tested integral of the tuple $v$ under
radius $r$, panel family $U$, and test function $t$. Write
$\mathcal T_0(v;t)$ for the corresponding reference test value. Hence
\[
 \mathcal T_{s_{n_d},U_d}(v;t)-\mathcal T_{0}(v;t)
 \longrightarrow0
\]
for every fixed tuple $v$ and test $t$. This is the tested approach property.
The argument is for each fixed test; it does not assert a uniform tail over
all tests or all label tuples.

The final chart isomorphisms follow grade by grade. For a fixed grade $d$,
the reference tests separate its finite-dimensional field: every nonzero
vector has a nonzero separated two-point evaluation. To obtain a smooth
test, keep the spectator insertion fixed and choose a separated pair where
the complex value is nonzero. Multiply the spectator by a unit complex
phase so that the real part of the value is positive. Continuity away from
the diagonal gives a small neighborhood on which that real part remains
positive. Integrating against a nonnegative smooth bump of integral one,
supported in that neighborhood, gives a smooth smeared test with nonzero
value. Thus smooth smeared reference tests also separate. Normalize to the
unit sphere. For
each point of that compact sphere, continuity gives a neighborhood on
which one such test has absolute value at least half its value at the
point. A finite subcover gives tests $t_1,\ldots,t_m$ and a constant
$\alpha>0$ such that every unit vector has
$\max_j|\mathcal T_0(v;t_j)|\geq\alpha$. The fixed finite set of
label charges and test suprema is bounded on that sphere. The tested
approach estimate therefore makes all corresponding errors less than
$\alpha/2$ for sufficiently large $k$. If the projected grade map killed a
unit vector, all its tested values would vanish, contradicting this lower
bound. Thus it is injective. The dimension bound from the selected table gives the reverse
dimension inequality; therefore the two dimensions agree and the map is
bijective. This tail may depend on $d$: there is no single finite index
after which every grade is simultaneously bijective. For the finitely many
indices before the grade-$d$ tail, choose any isomorphism using the same
dimension equality. This defines a scaling at every index. For each fixed
grade, the unit, reality, and signed symmetry identities hold on its
eventual isomorphism tail; the arbitrary isomorphisms used earlier need not
preserve those identities. The reference test-separation input here is the
all-grade free-detector result of
Section~\ref{sec:ex-lattice-all-grade-free-detectors}; the separate
low-grade trace calculation is not enough to justify separation in every
grade. The reviewed corrected-lift argument in
Section~\ref{sec:ex-lattice-trace-producer}
provides the compatible panel maps when required by the trace family.

\subsection{The gallery handoff}

The preceding argument proves the scale-selection, finite-cutoff extraction,
tested-approach, and fixed-grade eventual-isomorphism clauses. It does not
by itself prove the complete gallery statement. That handoff also needs the
physical-array receipts in the selected-run transfer. For the same
ultrafilter, hierarchy, and scaling these include the following. Let
$b_i>0$ be the selected arch scale at the $i$th postselected level, so
that $b_i\to\infty$; let $1$ and $\mathcal O$ be the unit and the
normalized curvature-square field of the reference grade-four table, and
$d_2$ the physical coefficient, as in
Section~\ref{sec:ex-lattice-trace-scribe-producer}. For a comparison table
$T$, write $\Theta_T$ for the scalar trace of the physical stress centered
by its one-point value in $T$, and $U_{i,4}$ for the grade-four chart from
the reference table to the selected table on its eventual isomorphism
tail. The receipts are:
convergence of the recentered physical stress to its prescribed principal
material tensor; for every comparison table that obeys the table hierarchy
laws and respects the selected table, real sequences $d_i,e_i$ such that,
eventually in $i$, the centered scalar trace satisfies
\[
 r_i^4U_{i,4}^{-1}\Theta_T=e_i\,1+d_i\,(4\mathcal O),
 \qquad b_id_i\longrightarrow d_2,\qquad b_ie_i\longrightarrow0.
\]
Thus $d_i,e_i\to0$ themselves. The final receipt is, for every comparison
table, convergence of $b_ir_i^4U_{i,4}^{-1}\Theta_T$ to
$4d_2\,\mathcal O$. The second direction
$4\mathcal O$ is the trace of the true-branch material tensor; the limit
of the zoomed observable is a multiple of $\mathcal O$, not of the unit. The arbitrary-null-
sequence case needs these same receipts after the selected scaling has
been constructed above. They are additional physical and field-array
estimates, not consequences of finite-dimensional compactness or of the
rank-one solve alone. The correction, one-dimensional
kernel, stable inverse, flux normalization, and physical/reference shell
comparison in Section~\ref{sec:ex-lattice-trace-producer} remain conditional
on their displayed analytic receipts.

In the selected public construction, the one-scale chart family comes from
the paid multi-scale request and its subsequent chart staging. This section
proves the scale-selection and tested-chart clauses of that assembly; the
physical-array, stress, and flux receipts needed for the full gallery remain
separate inputs.
